\chapter{Constraints}
\label{chap:constraints}

Constraints introduce kinematic relations between degrees of freedom that go beyond ordinary supports. A support prescribes selected degrees of freedom directly. A constraint instead relates the motion of different nodes or regions, for example by connecting a reference node to a surface, tying two independently meshed regions together, defining an idealized joint, or removing rigid-body motion from an otherwise free component.

The distinction between these constraint types is physical, not merely syntactic. A kinematic coupling imposes a rigid relation between selected motions. A structural/distributing coupling transfers generalized motion and loading between a reference point and a distributed region. A Tie creates a permanent compatibility relation between associated geometry. A Connector idealizes a joint between two points. RBM removes the rigid-body nullspace of a selected structural region. Contact, described at the end of this chapter, is unilateral and allows separation.

Constraints are used by the active analysis together with the selected support collectors. The numerical method used to enforce the resulting equations is controlled separately by \cmd{CONSTRAINTMETHOD}; see Chapter~\ref{chap:numerical-controls}.

\keywordpage{COUPLING}

\cmd{COUPLING} relates one master/reference node to a slave region. It is accepted at root and Assembly level after the model has been compiled. The reference is supplied by \key{MASTER}; \key{REF NODE} is accepted as an alias. The reference may name a node set containing exactly one node or a direct compiled node reference such as \val{42} or \val{INSTANCE.42}.

The coupled target is supplied by \key{SURFACE}. The legacy spellings \key{SLAVE} and \key{SFSET} remain aliases for the same target. FEMaster first resolves the name as a surface region and falls back to a node set. Thus element-based surfaces, node-based coupling regions, and existing native coupling decks use one common constraint path. If the same name exists in both namespaces, the surface region takes precedence.

Native FEMaster accepts two syntax styles. The legacy direct form places \key{TYPE=KINEMATIC} or \key{TYPE=STRUCTURAL} on \cmd{COUPLING} and follows it with exactly six selector values in the order
\[
(U_x,U_y,U_z,R_x,R_y,R_z).
\]
A component is active when the supplied value is greater than zero. Missing values are inactive. These six values are selectors rather than prescribed displacement or rotation magnitudes.

The scoped form is shared by the native and Abaqus readers. If \key{TYPE} is omitted, \cmd{COUPLING} opens a scope and exactly one child command selects its behavior: \cmd{KINEMATIC} or \cmd{DISTRIBUTING}. The child form uses inclusive Abaqus-style degree-of-freedom ranges. A line \val{first,last} activates all DOFs in that range, a single value activates only that DOF, and several lines may select several ranges. With no child data lines all six DOFs are selected.

\cmd{KINEMATIC} creates a kinematic coupling. The selected slave degrees of freedom follow the rigid kinematics induced by the reference node. This represents an idealized rigid flange, load-introduction point, or boundary where deformation of the coupled patch is suppressed in the selected directions.

\cmd{DISTRIBUTING} creates FEMaster's structural load-distribution coupling. The optional \key{COUPLING=CONTINUUM} and \key{COUPLING=STRUCTURAL} spellings are both accepted and deliberately map to the same internal \val{STRUCTURAL} coupling implementation. This transfers generalized loads and motions between the reference point and distributed region without making the complete target region rigid in the same sense as a kinematic coupling.

For a surface target, the distribution is based on the selected surface geometry. For a node-set target, the supplied node region itself defines the coupled region. A reference-point load combined with a distributing coupling represents a definite resultant force or moment independently of surface mesh density, whereas applying identical concentrated loads at every surface node changes the total load under mesh refinement.

The Abaqus reader supports the scoped \cmd{COUPLING}/\cmd{KINEMATIC}/\cmd{DISTRIBUTING} subset documented here. This does not imply support for every Abaqus coupling option; parameters outside the documented subset are not part of the FEMaster compatibility contract.

\begin{keywordparameters}
\key{MASTER} / \key{REF NODE} & required & Reference node or one-node node set. \\
\key{SURFACE} / \key{SLAVE} / \key{SFSET} & required & Nonempty slave surface or node set. \\
\key{TYPE} & optional & Native legacy direct form: \val{KINEMATIC} or \val{STRUCTURAL}. Omit to use a child scope. \\
\end{keywordparameters}

\begin{keyworddata}
Direct form & Six positive/inactive selectors for \(U_x,U_y,U_z,R_x,R_y,R_z\). \\
Scoped form & Data belong to \cmd{KINEMATIC} or \cmd{DISTRIBUTING} and use inclusive DOF ranges. \\
\end{keyworddata}

\exampleheading

\begin{dialectexamples}
\begin{femasterdeck}
** Legacy direct form
*COUPLING, MASTER=REF,
 SURFACE=LOADED_FACE, TYPE=STRUCTURAL
1., 1., 1., 0., 0., 0.

** Scoped form
*COUPLING, REF NODE=REF, SURFACE=LOADED_FACE
*DISTRIBUTING, COUPLING=STRUCTURAL
1, 3
\end{femasterdeck}
\begin{abaqusdeck}
*COUPLING, REF NODE=REF, SURFACE=LOADED_FACE
*DISTRIBUTING, COUPLING=STRUCTURAL
1, 3
\end{abaqusdeck}
\end{dialectexamples}

\keywordpage{KINEMATIC}

\cmd{KINEMATIC} is a child command of a \cmd{COUPLING} for which \key{TYPE} was omitted. It selects the kinematic coupling behavior and optionally restricts the active generalized degrees of freedom. The command is shared by the native and Abaqus readers.

Each data line contains an inclusive DOF range \val{first,last}. Valid DOF numbers are 1 through 6 and correspond to \(U_x,U_y,U_z,R_x,R_y,R_z\). If \val{last} is omitted, only \val{first} is selected. If all data lines are omitted, all six DOFs are active.

\begin{keyworddata}
\val{first, last} & Optional inclusive range of generalized DOFs 1--6. \\
\end{keyworddata}

\exampleheading
\begin{dialectexamples}
\begin{femasterdeck}
*COUPLING, MASTER=RP, SURFACE=FLANGE
*KINEMATIC
1, 3
6
\end{femasterdeck}
\begin{abaqusdeck}
*COUPLING, REF NODE=RP, SURFACE=FLANGE
*KINEMATIC
1, 3
6
\end{abaqusdeck}
\end{dialectexamples}

\keywordpage{DISTRIBUTING}

\cmd{DISTRIBUTING} is the structural/distributing child of a scoped \cmd{COUPLING} and is shared by the native and Abaqus readers. The optional \key{COUPLING} parameter accepts \val{CONTINUUM} or \val{STRUCTURAL}; both currently select FEMaster's structural load-distribution implementation.

Its optional DOF data uses the same inclusive 1--6 range syntax as \cmd{KINEMATIC}. With no range lines all generalized DOFs are selected.

\begin{keywordparameters}
\key{COUPLING} & optional & \val{CONTINUUM} (default) or \val{STRUCTURAL}; both map to FEMaster \val{STRUCTURAL}. \\
\end{keywordparameters}
\begin{keyworddata}
\val{first, last} & Optional inclusive range of generalized DOFs 1--6. \\
\end{keyworddata}

\exampleheading
\begin{dialectexamples}
\begin{femasterdeck}
*COUPLING, MASTER=RP, SURFACE=LOADED_FACE
*DISTRIBUTING, COUPLING=CONTINUUM
1, 3
\end{femasterdeck}
\begin{abaqusdeck}
*COUPLING, REF NODE=RP, SURFACE=LOADED_FACE
*DISTRIBUTING, COUPLING=CONTINUUM
1, 3
\end{abaqusdeck}
\end{dialectexamples}

\keywordpage{EQUATION}

\cmd{EQUATION} defines a homogeneous linear multi-point constraint from weighted nodal degrees of freedom. It is accepted at root or Assembly level after model topology has been compiled.

The first data line contains the total number of terms and must declare at least two. Following physical data lines contain one to four complete triples
\[
\text{target},\quad \text{dof},\quad \text{coefficient},
\]
until exactly the declared number of terms has been read. The degree of freedom is an integer from 1 through 6, corresponding to \(U_x,U_y,U_z,R_x,R_y,R_z\), and coefficients must be finite.

A target may be a direct node reference or a node set. If the first target is a single node, one scalar constraint equation is created. If the first target is an NSET, FEMaster expands the input into one equation per member of that set. Subsequent NSET targets are then paired with the first set in their existing compiled order and must have the same number of entries. Subsequent single-node targets are reused in every expanded equation.

For example, two equally sized sets \val{LEFT} and \val{RIGHT} can impose pairwise equality of one displacement component without manually writing one equation for every node pair. The ordering of the two sets therefore forms part of the equation definition.

\begin{keyworddata}
First line & Number of terms \(N\ge2\). \\
Following lines & One to four \val{target, dof, coefficient} triples per line until \(N\) terms are complete. \\
\end{keyworddata}

\exampleheading
\begin{dialectexamples}
\begin{femasterdeck}
*EQUATION
2
LEFT, 1, 1.0, RIGHT, 1, -1.0
\end{femasterdeck}
\begin{abaqusdeck}
*EQUATION
2
LEFT, 1, 1.0, RIGHT, 1, -1.0
\end{abaqusdeck}
\end{dialectexamples}

\keywordpage{TIE}

\cmd{TIE} creates a permanent kinematic connection between slave geometry and master geometry located within a prescribed search distance. The slave may be a node set or a surface set. The master may be a surface set or a line set. FEMaster selects the appropriate supported tie relation from these region types.

\key{DISTANCE} defines the geometric search tolerance used to associate slave locations with the master geometry. A distance below the interface separation leaves intended locations unmatched, whereas a distance extending to unrelated geometry admits additional candidates. \key{ADJUST=YES} permits the supported initial adjustment of slave geometry to the tie relation; the default is \val{NO}.

A Tie represents a bonded or otherwise permanently connected interface. Once the relation has been established, opening or sliding that violates the selected compatibility equations is not permitted. It therefore represents idealized bonded interfaces, mesh transitions, or connections that do not separate. \cmd{CONTACT} instead represents interfaces that may open, close, or lose contact.

The accuracy of a Tie depends on the geometric relation between the two discretizations. Highly nonmatching or strongly curved meshes and competing nearby master surfaces increase the sensitivity to the search distance.

Tie and contact solve different modelling problems. A Tie is an equality-type connection and remains active. Contact is unilateral: it transmits normal interaction only while active and permits subsequent separation. The stiffness of a tied interface is governed by the surrounding structure rather than an artificial penalty spring; physical compliance of an adhesive layer or joint is not represented by the Tie relation itself.

\begin{keywordparameters}
\key{MASTER} & required & Nonempty master surface or line set. \\
\key{SLAVE} & required & Nonempty slave node or surface set. \\
\key{DISTANCE} & required & Search/projection distance used to associate slave and master geometry. \\
\key{ADJUST} & optional & \val{NO} (default) or \val{YES}. \\
\end{keywordparameters}

\exampleheading

\begin{dialectexamples}
\begin{femasterdeck}
*TIE, MASTER=MASTER_FACE,
 SLAVE=SLAVE_NODES,
 DISTANCE=0.5, ADJUST=YES
\end{femasterdeck}
\begin{noabaqus}
Abaqus \cmd{TIE} is currently not imported by the FEMaster Abaqus reader, even though the native solver provides Tie constraints.
\end{noabaqus}
\end{dialectexamples}

\keywordpage{CONNECTOR}

\cmd{CONNECTOR} defines an idealized kinematic joint between two nodes. \key{NSET1} and \key{NSET2} each identify a node set containing exactly one node. \key{COORDINATESYSTEM} specifies the local connector axes, and \key{TYPE} selects which relative translations and rotations are constrained or released.

The supported connector types are \val{BEAM}, \val{HINGE}, \val{CYLINDRICAL}, \val{TRANSLATOR}, \val{JOIN}, and \val{JOINRX}. They represent different idealized joint kinematics. For example, a hinge is intended to release a rotational degree of freedom about its local hinge axis while constraining the remaining relative motion required by that joint idealization. A translator and a cylindrical connector likewise rely on their local axis system to define the permitted relative motion.

The coordinate system is therefore part of the physical definition, not a cosmetic orientation. Different local connector axes produce different joint kinematics even when the connector type and nodal positions remain unchanged.

Connectors idealize joints without resolving their internal geometry. Bearing deformation, bolt bending, local contact, joint clearance, and detailed connection stress require corresponding beam, shell, solid, spring, or contact representations.

A coupling can reduce a distributed region to a reference node before two reference nodes are connected. This construction makes the idealized joint definition independent of the surface mesh density.

\begin{keywordparameters}
\key{TYPE} & required & \val{BEAM}, \val{HINGE}, \val{CYLINDRICAL}, \val{TRANSLATOR}, \val{JOIN}, or \val{JOINRX}. \\
\key{NSET1} & required & First one-node set. \\
\key{NSET2} & required & Second one-node set. \\
\key{COORDINATESYSTEM} & required & Named coordinate system defining the connector axes. \\
\end{keywordparameters}

\exampleheading
\begin{dialectexamples}
\begin{femasterdeck}
*ORIENTATION, NAME=HINGE_AXES,
 TYPE=RECTANGULAR
1., 0., 0., 0., 1., 0.

*CONNECTOR, TYPE=HINGE,
 NSET1=PIN_A, NSET2=PIN_B,
 COORDINATESYSTEM=HINGE_AXES
\end{femasterdeck}
\begin{noabaqus}
Connector translation is currently not part of the supported Abaqus input subset.
\end{noabaqus}
\end{dialectexamples}

\keywordpage{RBM}

\cmd{RBM} suppresses the rigid-body modes of a selected structural element region without requiring the user to choose arbitrary fixed nodes. The target is specified by \key{ELSET} or its alias \key{SET}; if omitted, \val{EALL} is used.

A free three-dimensional body has six rigid-body modes: three translations and three rotations. These modes have no elastic strain energy and therefore produce a structural nullspace when no physical support removes them. RBM constructs constraint equations that remove the rigid-body motion of the selected region while avoiding the strong local bias that can arise from fixing one arbitrary corner node in several directions.

RBM removes the rigid-body nullspace without introducing a particular physical support point. It does not represent attachment through bearings, bolts, symmetry, contact, or other supports, and it does not complete missing physical boundary conditions.

RBM is also different from inertia relief. RBM changes the admissible kinematics by suppressing rigid motion. Inertia relief retains the free-body interpretation and balances the applied load with an equivalent rigid-body inertial field.

RBM addresses rigid-body motion itself. Other singularities can instead arise from disconnected parts, released mechanisms, missing Sections, unsupported rotational degrees of freedom, or incorrect constraint definitions and remain separate model conditions.

\begin{keywordparameters}
\key{ELSET} / \key{SET} & optional & Nonempty target structural element set; default \val{EALL}. \\
\end{keywordparameters}

\exampleheading
\begin{dialectexamples}
\begin{femasterdeck}
*RBM, ELSET=FREE_COMPONENT
\end{femasterdeck}
\begin{noabaqus}
No direct RBM-suppression keyword is currently imported by the supported Abaqus reader.
\end{noabaqus}
\end{dialectexamples}

\section{Contact}
\label{sec:contact}

Contact is a unilateral constraint between two surface regions. In contrast to a Tie, which permanently enforces compatible motion between associated regions, contact permits the surfaces to separate and transmits normal interaction only when the contact condition becomes active.

The current FEMaster contact formulation is frictionless dual-mortar surface-to-surface contact with an augmented-Lagrange normal-contact law. Contact acts over the physical overlap of slave and master facets and is nonlinear because the active contact region changes with deformation. Its response depends on the master/slave assignment, surface orientation, and penalty scale.

\keywordpage{CONTACT}

\cmd{CONTACT} defines one frictionless contact pair between a nonempty master surface set and a nonempty slave surface set. \key{MASTER} and \key{SLAVE} name the two regions. \key{PENALTY} provides the initial contact penalty stiffness. \key{CLEARANCE} shifts the zero-gap condition and defaults to zero. \key{FLIP=YES} reverses the master-surface normal used by the contact pair; the default is \val{NO}.

For each active slave contact degree of freedom, FEMaster uses an augmented-Lagrange relation of the form
\[
p_i=\max\left(0,\lambda_i-\varepsilon g_i\right),
\]
where \(g_i\) is the current normal gap, \(\lambda_i\) is the augmented-Lagrange multiplier, and \(\varepsilon\) is the effective penalty stiffness. With the sign convention used by the formulation, the \(\max\) operation suppresses tensile contact: surfaces may separate freely and normal contact pressure is generated only on the compressive side of the constraint.

The mortar formulation determines interaction from the current physical overlap of slave and master facets rather than simply tying a slave node to one permanently selected master node. This makes the contact response considerably less mesh-position dependent than a simple node-to-node contact law, but it also means that changing overlap topology can introduce non-smooth changes into a nonlinear equilibrium iteration.

The user-supplied \key{PENALTY} is a numerical stiffness scale, not a physical material constant. A small value permits greater penetration, whereas a large value can make the global tangent system poorly conditioned. The augmented-Lagrange procedure can increase the effective penalty when additional enforcement is required. The relevant scale is therefore set by the structural stiffness of the contacting bodies and evaluated through the resulting penetration and nonlinear conditioning.

\textbf{Master/slave assignment.} The mortar formulation is not algebraically symmetric with respect to the two labels: the dual interpolation and normal multipliers are associated with the slave side. Similar discretizations can admit either assignment, but changing it between comparable models also changes the discrete formulation. With strongly different meshes, assigning the finer or more detailed side as slave places the dual interpolation on that discretization; mesh sensitivity remains part of the numerical comparison.

\textbf{Normal orientation.} Contact sign depends on the surface orientation. A reversed surface orientation can make a physically separated pair appear active or prevent a penetrating pair from developing reaction. \key{FLIP} reverses the master normal for the contact pair but does not change the underlying element connectivity or surface side label.

\textbf{Penalty magnitude.} Penetration decreases as contact enforcement becomes stiffer, while excessive stiffness degrades the conditioning of the equilibrium iterations. The input number alone therefore does not define contact accuracy independently of geometry and structural stiffness.

\textbf{Analysis type.} Contact changes with deformation and belongs to nonlinear static analysis. A linear analysis cannot represent opening, closing, changing overlap, or an evolving unilateral active set.

\begin{keywordparameters}
\key{MASTER} & required & Nonempty master surface set. \\
\key{SLAVE} & required & Nonempty slave surface set. \\
\key{PENALTY} & required & Initial augmented-Lagrange penalty stiffness. \\
\key{CLEARANCE} & optional & Contact-clearance offset; default \val{0}. \\
\key{FLIP} & optional & \val{NO} (default) or \val{YES}; reverses the master-surface normal. \\
\end{keywordparameters}

\exampleheading
\begin{dialectexamples}
\begin{femasterdeck}
*SURFACE, NAME=MASTER_FACE
MASTER_ELEMENTS, S1
*SURFACE, NAME=SLAVE_FACE
SLAVE_ELEMENTS, S2

*CONTACT, MASTER=MASTER_FACE,
 SLAVE=SLAVE_FACE,
 PENALTY=1.0e8,
 CLEARANCE=0.0, FLIP=NO
\end{femasterdeck}
\begin{noabaqus}
The supported Abaqus reader currently does not import the Abaqus contact-definition family. Contact can be added with native FEMaster syntax using surface sets available in the model.
\end{noabaqus}
\end{dialectexamples}